%\documentclass[nofootinbib,floatfix,aps,pra,twocolumn]{revtex4}
%%  As of 17 Dec. 2002
%%  Last revised: SBT
%%  based on clas02b
%%  This is Sect VI ``The Proposed Research''
%\bibliographystyle{apsrev}


\documentclass[nofootinbib,floatfix,aps,pra,preprint]{revtex4}
%%%%%%%%%%%%%%%%%%%%%%%%%%%%%%%%%%%%%%%%%%%%%%%%%%%%%%%%%%%%%%%%%%%%%%%%%%%%%%%%%%%%%%%%%%%%%%%%%%%%%%%%%%%%%%%%%%%%%%%%%%%%
\usepackage{graphicx}

%TCIDATA{OutputFilter=LATEX.DLL}
%TCIDATA{Version=4.10.0.2347}
%TCIDATA{LastRevised=Wednesday, December 18, 2002 13:20:05}
%TCIDATA{<META NAME="GraphicsSave" CONTENT="32">}

\newcommand{\be}{\begin{equation}}
\newcommand{\ee}{\end{equation}}
\newcommand{\bea}{\begin{eqnarray}}
\newcommand{\eea}{\end{eqnarray}}
\topmargin -1.0cm
\oddsidemargin 0.cm 
\textheight 22.5cm 
\textwidth 17.0cm 
\sloppy

\input{tcilatex}

\begin{document}


% enumeration in items as A, B etc, not as 1, 2,...
\renewcommand{\labelenumi}{\Alph{enumi}.}

%\newpage
% sbt: Restart page counter at prescribed value
\setcounter{page}{1} \renewcommand{\baselinestretch}{1.1}

\textrm{\vspace*{-01.1cm} \noindent }

\begin{center}
\textrm{{\large \textsc{Systematic Construction of Better Spin-Density
Functionals}} \textsc{Abstract} }
\end{center}

\textrm{Practical implementation of a new method for systematic construction
of approximate density functionals is proposed. The construction builds the
functionals from a pre-selected kind of state (for example, determinant with
general spin orbitals or singles and doubles CI or AGP, etc.), so the
functionals will have known properties related directly to that type of
state. We will focus on spin topologies as they change with geometry because
of the difficulty these problems pose for many pesent-day DFT approximations
and because of their high relevance to catalysis, biochemical reactions, and
self-assembly of nanostructures. }

\textrm{The initial target will be on non-collinear spin configurations at
transition metal surfaces and ultrathin films, and related problems for
metal clusters. For rapidity of computation on complicated systems, we will
build functionals without explicit orbital dependence. Because the
functionals can have the properties of a known type of wave function built
in, the expectation is that many of the limitations of present-day
orbital-independent functionals will be overcome, so that a much greater
range of spin-dependent phenomena on chemically active surfaces can be
treated reliably and swiftly. }

\textrm{The method of construction involves systematic use of
transformations we have derived which map between and within elements of a
fiber bundle structure which we have proved is equivalent to the well-known
Levy-Lieb constrained search formulation of Density Functional Theory. }

\textrm{The proposal is to fund a postdoctoral associate, plus modest travel
for the co-author of the method (B.~Weiner) to be part of the project.
\newpage \noindent }

\begin{center}
\textrm{\textsc{Concept} }
\end{center}

\textrm{\setcounter{page}{1} \noindent \emph{Issues} }

\textrm{Density Functional Theory (DFT) has evolved to being among the major
quantum mechanical methods for predicting the properties of materials and
molecules from first principles. The 1998 Nobel Prize in Chemistry
recognized the remarkable feat accomplished via DFT. It avoids having to
cope explicitly with the mathematical complexity of the quantum mechanics of
hundreds to thousands of electrons treated explicitly. It does this by
converting the problem of predicting the properties of a molecule or
material into the study of a \emph{physically observable} function of three
spatial variables and a spin variable, the electron spin (number) density. }

\textrm{Brief surveys of the state of the art in DFT are found in Refs.~\cite%
{DFTrefs}. Particularly with the advent of so-called hybrid approximate
functionals (more below), there has been an outpouring of calculations on
molecular properties. The phrase ``approximate functionals'' points,
however, to the central intellectual challenge for predictive use of DFT.
The basic Hohenberg, Kohn, and Sham (HKS) theorems \cite{HKS} provide
assurance that the requisite mathematical objects, the spin-density
functionals, exist and have the needed properties. \emph{Unlike} other forms
of many-electron quantum mechanics however, the basic theorems give little
more than hints at the detailed properties of the exact functionals.
Approximations are needed therefore. But the challenge is stiffened, since
neither the theorems nor their proofs provide an explicit, systematic
framework for devising sequences of succesively more realistic approximate
functionals. The contrast with modern methods based explicitly on $N$%
-electron states (e.g. configuration interaction, coupled-cluster) is
striking. }

\textrm{As would be expected, this dilemma has generated a diversity of
responses: \vspace{-0.2cm} }

\begin{itemize}
\item \textrm{Use of empirical, numerical procedures to force the
approximate DFT functional to reproduce certain features and properties
determined from very-high-precision molecular wave-function (for example,
Ref.~\cite{Handy}); }

\item \textrm{Constraint of the approximate functional to have the known
mathematical properties of the ``exchange hole'' and ``correlation hole''
functions; pointers to the large original literature are in Sect. 3 of \cite%
{DFTrefs}(g); }

\item \textrm{Constraint of the approximate functional to satisfy
inequalities from coordinate scaling analysis \cite{scaling}; }

\item \textrm{Design of approximations that satisfy known asymptotic
constraints and can be calibrated to thermo-chemical data \cite{BeckeBaer}. }

\item \textrm{Use of explicit $N$-fermion theories terminated at some order
as implicit density functionals (``ab initio DFT'') \cite{RJBaiDFT} }
\end{itemize}

\textrm{It is important to distinguish non-empirical and empirical
functionals. Non-empirical functions incorporate no experimental data in
their parameterization (though many do incorporate information from
numerical studies, especially of the homogeneous electron gas). While
empirical approximations have enjoyed great popularity, they clearly cannot
be considered to be predictive in the full sense of the term. Since our goal
is the realization of the first-principles character of DFT, our attention
is directed to non-empirical functionals. }

\textrm{For all their successes, present-day non-empirical approximate
functionals developed these approaches (and variations and combinations)
have several well-known shortcomings when used on real materials and
molecules. The limitations pertinent to this project all involve the
interplay of spin configurations, system energetics and geometry, and
mechanical properties. Critical limitations include }

\begin{itemize}
\item \textrm{\vspace{-0.21cm} }

\item \textrm{The dilemma of choice between unphysical spin-symmetry
breaking and unphysically high energies as systems are separated (diluted) 
\cite{PerdewSavin}; \vspace{-0.21cm} }

\item \textrm{The incomplete treatment of general spin density topologies
(e.g.~non-collinear magnetism), typically addressed either by ingenious
individual mapping to a presumably equivalent problem (``unwinding'')
amenable to a standard DFT approximation \cite{Kleinman98,Oda} or else by
theorems that generalize standard DFT approximations to handle a particular
spin density topology \cite{Capelle00}; \vspace{-0.21cm} }

\item \textrm{A poor description of ground states that are strong admixtures
of nearly-degenerate single-determinant configurations (in quantum chemistry
terms, a static correlation problem); \vspace{-0.21cm} }

\item \textrm{The failure to bind negative ions correctly (or at all in some
cases) \cite{NRSBT97} and improper behavior of the ground state energies of
isoelectronic sequences of positive atomic ions \cite{Davidson98}; \vspace{%
-0.21cm} }

\item \textrm{An erratic balance (from system to system) in quality of
calculated cohesive energies (atomization energies) vs. lattice constants
(bond lengths) vs. bulk moduli (ground state vibrational frequency) \cite%
{Zupan98,SBT97a}. \vspace{-0.21cm} }
\end{itemize}

\textrm{These limitations are particularly important for the prediction of
interactions with open-shell systems (clusters, nanostructures, surfaces),
interactions at metallic surfaces or at sites with metal atoms, and for
tuning surface, cluster, or nanostructure behavior by tailoring of spin
populations. These physical and chemical aspects are relevant, in turn, to
issues of structure, mechanical stability, and efficacy of heterogeneous
catalysts and self-assembled bio-nano materials, to give but two examples. }

\textrm{\vspace{1cm} \noindent \emph{New Approach} }

\textrm{These limitations on the chemical and physical reliability of
present-day non-empirical functionals can be traced to the details of how
those functionals handle spin configurations. In a certain sense the point
is obvious: because DFT parameterizes the $N$-electron quantum mechanical
ground state via the electron spin number density, the details of the
interplay between anti-symmetrization and spatial distribution in the ground
state are concealed. }

\textrm{Though the problem is obvious, that does not mean that a systematic
solution also is obvious. Our approach is to devise, then implement an
explicit, systematic characterization of the mappings from the known
spin-configuration properties of a specified kind of state function to a
functional. This proposal is to fund the initial \emph{implementation}
phase; we have completed the formal mathematical analysis already. By use of
this constructive DFT formalism already developed by B.~Weiner and me, we
will devise and test new functionals as well as improve existing ones.
Weiner is on the Physics faculty at Penn State, with teaching
responsibilities at the Dubois undergraduate campus. He has visited QTP on
several occasions. He will be a partner in this project, as detailed below. }

\textrm{The constructive DFT formalism is the product of our recent work on
formal analytical properties of DFT. The mathematical foundations of the
approach are published \cite{BLWSBT97,BLWSBT98,BLWSBT99}. To preserve
competitive advantage (more discussion below), we have not published the
explicit procedure for constructing functionals. A summary of the
time-independent case follows. The time-dependent analogue is substantially
more intricate and involves a different variational principle. Nevertheless,
the time-dependent and time-independent methodologies work in essentially
the same way so far as constructing the functionals are concerned. }

\textrm{First, recall that the Hohenberg-Kohn-Sham (HKS) theorems are: %
\renewcommand{\labelenumi}{\Roman{enumi}.} }

\begin{enumerate}
\item \textrm{\vspace{-0.21cm} }

\item \textrm{The ground state energy of a many-electron system is a
functional of the electron spin number density alone. It consists of a
universal functional (i.e., a functional independent of the external
potential) of the electron spin number density plus the energy from the
total external potential $E_{tot}[n_{0}] = F[n_{0}] + \int {}{} \, v_{ext}({}%
) \, n_{0}({}) $ \vspace{-0.21cm} }

\item \textrm{The total energy has a variational minimum with respect to the
electron spin density at the exact ground state spin density: $%
E_{tot}[n_{0}] = \min_{n} E_{tot}[n] \;\; \forall \;\; n({}) \, \ni \, \int {%
}{}\, n({}) = N $ \vspace{-0.21cm} }

\item \textrm{The variational minimization problem can be set into
correspondence with a one-electron problem for the spin densities of an
auxiliary system of non-interacting Fermions (the Kohn-Sham system) whose
spin-orbitals are eigenfunctions of a purely local effective Hamiltonian. 
\vspace{-0.21cm} }
\end{enumerate}

\textrm{\renewcommand{\labelenumi}{\Alph{enumi}.} Today these statements
usually are proved via the constrained search approach due to Levy \cite%
{LevySem2,Levy79,LevyPerdew85} and Lieb \cite{Lieb}. Motivated originally by
spin-symmetry-breaking problems in DFT, the Weiner-Trickey reformulation
makes explicit the internal structure of constrained search as well as
resolving some critical, subtle difficulties \cite%
{BLWSBT97,BLWSBT98,BLWSBT99}. The proposed research depends primarily on the
outcomes of those papers and only secondarily on the lengthy, formal
analysis that they present. Thus a summary of the major technical components
will suffice to outline the underpinnings and show why it is indeed a
constructive methodology. Those main components include:\vspace{-0.07cm} }

\begin{enumerate}
\item \textrm{\vspace{-0.15cm} }

\item \textrm{Formulation as a linear problem in $N$-electron density
operators and their contraction maps to the spin density (rather than as a
bi-linear problem in Hilbert space states); }

\item \textrm{Classification of all $N$-electron density operators (pure 
\emph{and} ensemble) according to equivalence classes for each electron spin
density into a well-characterized fiber bundle structure;\vspace{-0.21cm} }

\item \textrm{Establishment of the existence of a Hohenberg-Kohn universal
variational energy functional on that bundle;\vspace{-0.21cm} }

\item \textrm{Proof of the existence of required functional derivatives of
the HK functional on that bundle (and proof that such derivatives may not
exist elsewhere);\vspace{-0.21cm} }

\item \textrm{Characterization of the group of invariances of the kernel of
the contraction map from the $N$-electron density operator to the
one-electron spin density;\vspace{-0.21cm} }

\item \textrm{Use of the properties of that invariance group to derive
explicit transformations for two cases }

\begin{itemize}
\item \textrm{transformations that leave the pure-state segment of a given
fiber invariant (i.e. move among states \emph{intra}-fiber), }

\item \textrm{transformations that go from the pure-state segment of one
fiber to the pure-state segment of another (i.e. move among states \emph{%
inter}-fiber) }
\end{itemize}

\textrm{and analogously for the set of fiber segments corresponding to a
specified ensemble. \vspace{-0.21cm} }
\end{enumerate}

\textrm{Though we did not initially recognize it, the closest antecedent to
our approach probably is the work of Englisch and Englisch \cite{Englisch}
and the slightly prior work of Zumbach and Maschke \cite{Zumbach}. Those
authors focused on the construction of all the states that contract to a
given density, on the existence of functional derivatives needed to carry
out the variation principle, etc. Both pairs of authors in fact gave
constructions for the wave functions that map to a specified density. The
constructions are quite difficult to use in constructing practical
functionals however, because in essence the constructions are a full
Configuration Interaction in a particularly cleverly devised basis. }

\textrm{Our approach differs most significantly in focusing on the linear
nature of the problem when posed in terms of density operators and
therefore, in the constructive part, items E and F in the foregoing list.
The linear nature of the formulation allows us to elucidate a group
structure and group parameterization of the transformations among states.
This is the core of our constructive approach. We can select a well-known
class of approximate wavefunctions (for example, single CI, or singles and
doubles CI, or AGP, etc. from quantum chemistry) and \emph{construct} the
form of functional which corresponds to that class. The supporting logic can
be summarized as follows:\vspace{-0.21cm} }

\begin{itemize}
\item \textrm{A given type of approximate state corresponds to a particular
sector of each fiber;\vspace{-0.21cm} }

\item \textrm{A functional on that type of state consists of a path joining
those selected segments such that the elements connected are the energy
minima for each such segment \vspace{-0.21cm} }

\item \textrm{That path can be constructed explicitly in terms of the group
parameters for the invariances of those segments. }
\end{itemize}

\textrm{Because rapidity of computation is essential for the predictive
treatment of complicated systems (\emph{e.g.} the quantum mechanical domain
in a multi-scale simulation), we will build functionals without explicit
orbital dependence. Though not the most fashionable approach, particularly
in the Chemistry part of the DFT community, it is in keeping with the
fundamental proposition of DFT, to wit, that the density (and not some
orbitals that generate it) is the fundamental quantity with which to
parameterize and $N$-Fermion ground state. }

\textrm{Because single determinant reference functions are critical in the
conventional implementation of HKS DFT, and because spin symmetry problems
are ubiquitous in such implementations, a particularly appealing route is to
study the density operators corresponding to Generalized Hartree Fock wave
functions. Since we have already shown how the Fukutome symmetry classes 
\cite{Fukutome} for such determinants exhaust all possible Kohn-Sham states
so far as spin topologies are concerned, the connection should be extremely
fruitful. }

\textrm{An obvious area of relevance for this approach is systems with
non-collinear spin densities. From the large literature on such systems we
pick but one example, the ``unwinding'' technique used by Bylander and
Kleinman \cite{Kleinman98} to treat fcc Fe. That technique is equivalent to
a special case mapping onto the seventh Fukutome class (``torsional spin
density waves'') followed by treatment of the component KS equations via the
ordinary local spin density approximation (LSDA). See also the closely
related paper by Oda et al.\cite{Oda} Also the theorems proved by Capelle
and Oliveira \cite{Capelle00} for non-collinear magnetism are particular
realizations of our general results regarding Fukutome classes. Recently
Kleinman \cite{Kleinman99} has put forth a functional specifically for
spiral spin density waves. Another application of our approach will be to
derive that functional from the antecedent density operator class. Note that
this general idea has some roots also in the work of von Barth and Hedin 
\cite{vBH72}. Our contribution is to reformulate the theory to give a
procedure for \emph{building } functionals that behave properly in the sense
of reproducing the known behavior of a particular kind of state. }

\textrm{It appears that analogous reasoning should hold for current density
functional theory (CDFT; references in \cite{VignaleRasolt}) but we have not
explored that supposition thus far. Doing so would leverage my present NSF
funding (DMR-0218957), which is under the Information Technology for
Research program in Division of Materials Research. A major component of
that award is to implement CDFT in the Boettger-Trickey computer code GTOFF.
Simultaneous formal investigation of CDFT would be immediately implementable
and testable computationally therefore. }

\begin{center}
\textrm{\textsc{Opportunity and Plan} }
\end{center}

\textrm{The scientific opportunity is to demonstrate the construction of
better functionals by use of our constructive formal DFT methods. The route
forward is clear, yet we have not been able to exploit it for lack of
personnel. In principle there also is a funding opportunity, except for an
impasse that we propose to resolve by this project. The impasse is that
proposals to both DOE and NSF have been rejected on the grounds that we have
not so far constructed a new functional. We have not done so because we need
more human resources. The mathematical formalism has been developed without
support as a kind of hobby. Such a marginal effort cannot support the
development and stringent testing of new functionals. So the impasse arises:
no funding implies no new functionals but no new functionals implies no
funding. }

\textrm{To escape the impasse we have done and/or are doing the following: }

\textrm{A. I have submitted an internal proposal to the College of Liberal
Arts and Sciences (``Cooperative Project on Improving Density Functionals'',
25 October 2002). Included in it is a commmitment to the remaining steps
listed here. The internal proposal is to jump-start the effort by short-term
(approximately 5.5 months) funding of a post-doctoral associate who is
well-prepared in DFT to work on functionals with us. At least two people
have been identified who are strong candidates, without even trying to
recruit such a person. That internal proposal is pending. }

\textrm{B. The present proposal to PRF is for funding of that same
post-doctoral for two years as well as modest travel and incidentals
primarily so that Weiner can work at Florida in the summer with us. Together
with step A, we believe that this would provide the funding base to produce,
test, and publish perhaps two new functionals and two to three improvements
of existing functionals. }

\textrm{C. Weiner will be submitting a Research at Undergraduate
Institutions proposal to NSF-DMR to work on Current Density Functional
Theory as it relates to the implementation phase in my NSF ITR grant; recall
previous remarks about CDFT. In both cases his expertise will be focused on
better functionals, so there is synergy. }

\textrm{D. A new (Fall 2002) Chemistry graduate student, Joe Stember
approached me in September, 2002, and asked to study formal DFT with the
intent of doing his research in that area. I have started tutoring him with
notes from my PHY 7097 course, taught first in Fall 2001. Stember has a
Grinter fellowship from Univ. of Florida, which gives full support for 4
years. Though he has not selected a thesis advisor at this writing, I am the
only faculty member with whom he is working as an individual. If he does
become my student, the obvious direction for his Ph.D.~thesis will be this
project. }

\textrm{The plan is to begin by constructing functionals explicitly for
non-collinear spin topologies from the corresponding Fukutome class. We will
test these on ``toy'' systems using the Boettger-Trickey GTOFF code. An
example is the H mononlayer. Though physically unrealistic, it poses a
difficult problem with spin configurations as the lattice parameter is
enlarged. Thus it is a good shake-down test for comparison with published
functionals. We then will consider systems such as discussed above in
respect to the work of Kleinman and co-workers. In particular, we will test
for the dependence of spin topology on lattice constant. We also will
compare with the model Hamiltonian treatment of Uzdin \emph{et al.} \cite%
{Uzdin99} and the DFT studies by Kurz \emph{et al.} \cite{Kurz00} of the
non-collinear magnetism of transition-metal ultra thin films. Issues to be
assessed include the physical realism of the resulting spin topologies, the
extent to which the new functional gives qualitatively and/or quantitatively
different descriptions of the surface chemistry, etc. A related issue is the
strain-dependent cusping of the spin-magnetic moments of the Fe monolayer
found originally by Boettger and extended to all 3d transition metal
monolayers by the two of us \cite{JCB93}. The question is to what extent
this phenomenon, which could be useful in tuning surface reactivity, is a
sharp consequence of the relatively simple DFT approximation we used (LSDA). 
}

\textrm{If these explorations are as successful as we expect (or even
moderately encouraging), we will implement the new functional at
prototype-level coding in the ACES-II program package \cite{ACESII} and test
for performance on the notoriously difficult case of Cr$_2$ as the bond
length is extended. }

\textrm{Second, we will examine popular spin-dependent functionals based on
scaling and inequalities, \emph{e.g.} PBE \cite{PBE}, to see how they can be
``reverse-engineered'' to uncover their connection with a particular class
of wave function or ensemble. Characterizing such a connection would enable
one to state with precision and clarity the limits of predictive utility for
that particular functional. Additionally, this reverse-engineering should
enable us to propose refinements and/or extensions of the existing
functionals to improve the fidelity with which they track to a known class
of wave function or ensemble. }

\textrm{Assuming, as expected, that the new or modified functionals will
achieve substantially greater physical realism as well as more generality
than existing ones, we may also attempt to build a functional that follows
from an explicitly correlated wave function, probably a geminal-based one,
such as AGP (which has been used to provide a formal basis for the structure
of electron propagator approximations) \cite{AGP}. The objective of this
construction would be to assess the way in which correlation contributions
are handled by the functional. }

\textrm{However, if the pursuit of functionals from other Fukutome classes
(which themselves are correlated in the sense of being multi-determinantal,
since they use general spin orbitals) seems more promising, we shall follow
that route. Since we have lacked the opportunity until now to try out our
constructive approach, detailed research plans beyond those just enunciated
are more than a bit speculative. }

\textrm{\vspace{1.1cm} }

\begin{thebibliography}{99}
\bibitem{} \textrm{%\setcounter{page}{1}
}

\bibitem{DFTrefs} \textrm{{\small (a) ``The Density Functional Formalism,
its Applications and Prospects'', R.O.~Jones and O.~Gunnarsson, Rev. Mod.
Phys. \textbf{61}, 689-746 (1989); (b) R.G.~Parr and W.~Yang, \textit{%
Density Functional Theory of Atoms and Molecules} (Oxford, NY, 1989); (c)
R.M.~Dreizler and E.K.U.~Gross, \textit{Density Functional Theory}
(Springer, Berlin, 1990); (d) E.S.~Kryachko and E.V.~Lude\~na, \textit{%
Energy Density Functional Theory of Many-Electron Systems} (Kluwer,
Dordrecht, 1990); (e) ``Density Functional Theory of Electronic Structure'',
W.~Kohn, A.D.~Becke, and R.G.~Parr, J. Phys. Chem. \textbf{100}, 12974-80
(1996); (f) H.~Eschrig, \textit{The Fundamentals of Density Functional Theory%
}, Teubner Texte f\"ur Physik \textbf{32} (Teubner, Stuttgart and Leipzig,
1996); (g) ``A Critical Assessment of Density Functional Theory with Regard
to Applications in Organometallic Chemistry'', A.~G{\"o}rling, S.B.~Trickey,
P.~Gisdakis, and N.~R{\"o}sch, in \textit{Topics in Organometallic Chemistry}%
, vol.\ 4, P.~Hoffmann and J.M.~Brown eds. (Springer, Berlin, 1999) 109-63. }%
}

\bibitem{HKS} \textrm{{\small (a) ``Inhomogeneous Electron Gas'',
P.~Hohenberg and W.~Kohn, Phys. Rev. \textbf{136}, B864-71 (1964); (b)
``Self-Consistent Equations Including Exchange and Correlation Effects'',
W.~Kohn and L.J.~Sham, Phys. Rev. \textbf{140}, A1133-38 (1965). }}

\bibitem{Handy} \textrm{{\small ``Learning from Exchange-Correlation
Potentials'', W.H.~Green, D.J.~Tozer, and N.C.~Handy, Chem. Phys. Lett. 
\textbf{290}, 465-72 (1998) and references therein. }}

\bibitem{scaling} \textrm{{\small (a) M.~Levy and A.~G\"orling, Phys. Rev. A 
\textbf{51}, 2851-55 (1995); (b) C.~Filippi, X.~Gonze, and C.J.~Umrigar, in 
\emph{Recent Developments and Applications of Modern Density Functional
Theory}, J.A.~Seminario ed. (Elsevier, Amsterdam, 1996), 295-326. }}

\bibitem{BeckeBaer} \textrm{{\small (a) ``Density-functional Exchange-energy
Approximation with Correct Asymptotic Behavior'', A.D.~Becke, Phys. Rev. A 
\textbf{38}, 3098-3100 (1988); (b) ``Exchange-correlation Potential with
Correct Asymptotic Behavior'', R.~van Leeuwen and E.J.~Baerends, Phys. Rev.
A \textbf{49}, 2421-31 (1994). }}

\bibitem{RJBaiDFT} \textrm{{\small ``\emph{Ab Initio} Density Functional
Theory: OEP-MBPT(2). A New Orbital-dependent Correlation Functional'',
I.~Grabowski, S.~Hirata, S.~Ivanov, and R.J.~Bartlett, J. Chem. Phys. 
\textbf{116}, 4415-25 (2002). }}

\bibitem{PerdewSavin} \textrm{{\small ``Escaping the Symmetry Dilemma
through a Pair-Density Interpretation of Spin-Density Functional Theory'',
J.P.~Perdew, A.~Savin, and K.~Burke, Phys. Rev. A \textbf{51}, 4531-41
(1995). }}

\bibitem{Kleinman98} \textrm{{\small ``Full Potential \textit{ab initio}
Calculations of Spiral Spin Density Waves in fcc Fe'', D.M.~Bylander and
L.~Kleinman, Phys. Rev. B \textbf{58}, 9207-11 (1998). }}

\bibitem{Oda} \textrm{{\small ``Fully Unconstrained Approach to Noncollinear
Magnetism: Application to Small Fe Clusters'', T.~Oda, A.~Pasquarello, and
R.~Car, Phys. Rev. Lett. \textbf{80}, 3622-25 (1998). }}

\bibitem{Capelle00} \textrm{{\small ``Density Functional Theory for
Spin-density Waves and Antiferromagnetic Systems'', K.~Capelle and
L.N.~Oliveira Phys. Rev. B \textbf{61}, 15228-40 (2000). }}

\bibitem{NRSBT97} \textrm{{\small ``Negative Ions in Density Functional
Theory - Comment on J. Chem. Phys. \textbf{105}, 862 (1996)'', N. R{\"o}sch
and S.B. Trickey, J. Chem. Phys. \textbf{106}, 8940-41 (1997) and references
therein. }}

\bibitem{Davidson98} \textrm{{\small ``Ground-state Energies of
Isoelectronic Atomic Series from Density-Functional Theory: Exploring the
Accuracy of Density Functionals'', A.A.~Jarzecki and E.R.~Davidson, Phys.
Rev. A \textbf{58}, 1902-08 (1998). }}

\bibitem{Zupan98} \textrm{{\small ``Pressure-induced Phase Transition in
Solid Si, $\mathrm{SiO_{2}}$, and Fe: Performance of LSD and GGA Density
Functionals'', A.~Zupan, P.~Blaha, K.~Schwarz, and J.P.~Perdew, Phys. Rev. B 
\textbf{58}, 11266-272 (1998). }}

\bibitem{SBT97a} \textrm{{\small ``Benchmark Comparison of Gradient
Dependent and Local Density Calculations for Bulk Si and Al'', S.B.~Trickey,
Internat. J. Quantum Chem. \textbf{61}, 641-46 (1997). }}

\bibitem{LevySem2} \textrm{{\small ``Coordinate Scaling Requirements for
Approximating Exchange and Correlation'', M.~Levy, in \emph{Recent
Developments and Applications of Modern Density Functional Theory},
J.A.~Seminario ed. (Elsevier, Amsterdam, 1996), 3-24. }}

\bibitem{Levy79} \textrm{{\small ``Universal Variational Functionals of
Electron Densities, First-Order Density Matrices, and Natural Spin Orbitals
and Solution of the v-Representability Problem'', M.~Levy, Proc. Natl. Acad.
Sci. USA \textbf{76}, 6062-65 (1979). }}

\bibitem{LevyPerdew85} \textrm{{\small ``The Constrained Search Formulation
of Density Functional Theory'', M.~Levy and J.P.~Perdew in \emph{Density
Functional Methods in Physics}, R.M.~Dreizler and J.~da Providencia (Plenum,
NY, 1985) 11-30. }}

\bibitem{Lieb} \textrm{{\small ``Density Functionals for Coulomb Systems'',
E.H. Lieb, Internat. J. Quantum. Chem. \textbf{24}, 243-77 (1983). }}

\bibitem{BLWSBT97} \textrm{{\small ``Fukutome Symmetry Classification of the
Kohn-Sham Auxiliary One-Matrix and its Associated State or Ensemble'',
B.~Weiner and S.B.~Trickey, Internat. J. Quantum Chem. \textbf{69} 451-60
(1998) Also at http://www.personal.psu.edu/bqw/list\_pap.html }}

\bibitem{BLWSBT98} \textrm{{\small ``Time Dependent Variational Principle in
Density Functional Theory'', B.~Weiner and S.B.~Trickey, Adv. Quantum
Chemistry (Academic, San Diego), \textbf{35}, 217-48 (1999). Version with a
shortened introduction is at http://www.personal.psu.edu/bqw/list\_pap.html }%
}

\bibitem{BLWSBT99} \textrm{{\small ``State Energy Functionals and
Variational Equations in Density Functional Theory'', B.~Weiner and
S.B.~Trickey, J. Molec. Struc. (Theochem) \textbf{501}-\textbf{02}, 65-83
(2000). Also at http://www.personal.psu.edu/bqw/list\_pap.html }}

\bibitem{Englisch} \textrm{{\small (a) ``Exact Density Functionals for
Ground-State Energies I. General Results'', H.~Englisch and R.~Englisch,
Phys. Stat. Solidi (b) \textbf{123}, 711-21 (1984) (b) ``Exact Density
Functionals for Ground-State Energies II. Details and Remarks'', H.~Englisch
and R.~Englisch, Phys. Stat. Solidi (b) \textbf{124}, 373-79 (1984). }}

\bibitem{Zumbach} \textrm{{\small ``New Approach to the Calculation of
Density Functionals,'' G.~Zumbach and K.~Maschke, Phys. Rev. A \textbf{28},
544-54 (1983); erratum Phys. Rev. A \textbf{29}, 1585-86 (1984). }}

\bibitem{Fukutome} \textrm{{\small ``Unrestricted Hartree-Fock Theory and
Its Applications to Molecules and Chemical Reactions'', H.~Fukutome,
Internat. J. Quantum Chem. \textbf{20}, 955-1065 (1981). }}

\bibitem{VignaleRasolt} \textrm{{\small ``Magnetic Fields and Density
Functional Theory'', G.~Vignale, M.~Rasolt, and D.J.W.~Geldart in Adv.
Quantum Chem. \textbf{21}, 235-53 (1990). }}

\bibitem{Kleinman99} \textrm{{\small ``Density Functional for Noncollinear
Magnetic Systems'', L.~Kleinman, Phys. Rev. B \textbf{59}, 3314-17 (1998). }}

\bibitem{vBH72} \textrm{{\small ``A Local Exchange-Correlation Potential for
the Spin Polarized Case: I'', U.~von Barth and L.~Hedin, J. Phys. C: Sol.
St. Phys. \textbf{5}, 1629-42 (1972). }}

\bibitem{Uzdin99} \textrm{{\small ``Noncollinear Magnetism of Fe/CR Films
and Multilayers'', V.M.~Uzdin, N.S.~Yarseva, and S.V.~Yartsev, J. Mag. Mag.
Matl. \textbf{196-97}, 70-72 (1999). }}

\bibitem{Kurz00} \textrm{{\small ``Noncollinear Magnetism of Cr and Mn
Monolayers on Cu(111)'', Ph.~Kurz, G.~Bihlmayer, and S.~Bl\"ugel, J. Appl.
Phys. \textbf{87} 6101-03 (2000). }}

\bibitem{JCB93} \textrm{{\small ``Theoretical Strain-dependent Properties of
Square and Hexagonal Fe Monolayers'', J.C.~Boettger, Phys. Rev. B \textbf{48}%
, 10247-53 (1993); ``First-Principles Systematics of Ordered Metallic
Monolayers: 3D Transition Metals'', S.B.~Trickey and J.C.~Boettger, Invited
lecture, Fifth Chemical Congress of North American, Canc\'un, M\'exico,
11-15 Nov. 1997 (unpublished). }}

\bibitem{ACESII} \textrm{{\small ACES II is a program product of the Quantum
Theory Project, University of Florida. Authors: J.F. Stanton, J. Gauss, S.
A. Perera, J.D. Watts, M. Nooijen, Anthony Yau, N. Oliphant, P. G. Szalay,
W.J. lauderdale, S.R. Gwaltney, S. Beck, A. Balkov{\'a} D. E. Bernholdt,
K.-K Baeck, P. Rozyczko, H. Sekino, C. Huber, Ji\u{r}i{\'a} Pittner and R.J.
Bartlett. Integral packages included are VMOL (J. Alml{\"o}f and P.R.
Taylor); VPROPS (P. Taylor) ABACUS; (T. Helgaker, H.J. Aa. Jensen, P. J{\o }%
rgensen, J. Olsen, and P.R. Taylor). }}

\bibitem{PBE} \textrm{{\small ``Generalized Gradient Approximation Made
Simple'', J.P.~Perdew, K.~Burke, and M.~Ernzerhof, Phys. Rev. Lett. \textbf{%
77}, 3865-68 (1997), erratum Phys. Rev. Lett. \textbf{78}, 1396 (1997) }}

\bibitem{AGP} \textrm{{\small ``The AGP wavefunction and Its Relation to
Other Descriptions of Electronic Structure'', J.V.~Ortiz, B.~Weiner, and
Y.~\"Ohrn, Internat. J. Quantum Chem. S-\textbf{15}, 113-28 (1981) }}
\end{thebibliography}

\end{document}
